% Neuron: 150 words summary
% PNAS 250 words abstract + 50-120 words significant statement
% current: 217 words

Memory consolidation is the process by which temporary, malleable memories are transformed into more stable, longer-lasting forms. On a coarse anatomical scale, consolidation redistributes memories in the brain, but it remains poorly understood how these changes manifest themselves on the finer, cellular scale of neuronal engrams and how they relate to the cognitive level. 
In this study, we developed a phenomenological model of engram dynamics under systems consolidation. %In our model, changes in neural representations are driven by reactivation, and we capture how engrams evolve over long timescales. 
The model describes consolidation as a brain-wide phenomenon, where memories deterministically follow a trajectory through a space of patterns distributed among brain regions. It captures a broad range of features of memory consolidation, including selective consolidation, semantization, and power-law forgetting. 
In the model, consolidation is accompanied by population-level changes in neuronal representations that resemble the widely observed phenomenon of representational drift. 
%Recent experimental work has shown that the mapping of task variables to population activity changes over days to weeks, even in the absence of task learning. Our model suggests that this representational drift may be a consequence of memory consolidation. 
When only a subset of neurons is observed, the deterministic dynamics of the model can appear stochastic, and a readout of task features deteriorates over time even when a stable readout exists for the full system. Our model offers a dynamical systems perspective on memory consolidation as a distributed process, moving beyond the classic region-centered view, and provides a functional interpretation of drift as a means of redistributing engrams for improved memory retention.

%Memory consolidation is the process by which temporary, malleable memories are transformed into more stable, longer-lasting forms. On a coarse anatomical scale, consolidation redistributes memories in the brain, but it remains poorly understood how these changes manifest themselves on the finer, cellular scale of neuronal engrams and how they relate to the cognitive level.  In this study, we developed a phenomenological model of engram dynamics under systems consolidation.   The model describes consolidation as a brain-wide phenomenon, where memories deterministically follow a trajectory through a space of patterns distributed among brain regions. It captures a broad range of features of memory consolidation, including selective consolidation, semantization, and power-law forgetting.  In the model, consolidation is accompanied by population-level changes in neuronal representations that resemble the widely observed phenomenon of representational drift.  When only a subset of neurons is observed, the deterministic dynamics of the model can appear stochastic, and a readout of task features deteriorates over time even when a stable readout exists for the full system. Our model offers a dynamical systems perspective on memory consolidation as a distributed process, moving beyond the classic region-centered view, and provides a functional interpretation of drift as a means of redistributing engrams for improved memory retention.